\chapter{Flusso operativo}
\label{chap:lavorazione}

\section{Finalità e ambito}

La pipeline ricostruisce un grafo a partire da fotografie multi-vista dello
stesso nodo idraulico. Individua i simboli Shurikode, associa le detection
riferite allo stesso oggetto fisico e, quando disponibile e abilitata, usa la
geometria prodotta da VGGT come supporto per la riconciliazione delle istanze e
la stima delle connessioni. Il risultato è un grafo JSON corredato da artefatti
per il controllo e la valutazione.

\begin{figure}[H]
  \centering
  \begin{tikzpicture}[
  node distance=6mm and 7mm,
  opstep/.style={draw=ManualBlue, rounded corners, fill=ManualLight,
               align=center, minimum height=13mm, text width=29mm,
               font=\small},
  result/.style={draw=ManualGreen, rounded corners, fill=white,
                 align=center, minimum height=13mm, text width=31mm,
                 font=\small},
  arrow/.style={-{Latex[length=2mm]}, thick, draw=ManualDark}
]
\node[opstep] (data) {\textbf{1.} Immagini\\multi-vista};
\node[opstep, right=of data] (env) {\textbf{2.} Configurazione\\della lavorazione};
\node[opstep, right=of env] (run) {\textbf{3.} Esecuzione\\della pipeline};
\node[result, below=of run] (out) {\textbf{4.} Grafo\\ricostruito};
\node[opstep, left=of out] (check) {\textbf{5.} Verifica e\\valutazione};
\node[result, left=of check] (deliver) {\textbf{6.} Output per\\utilizzo e consegna};

\draw[arrow] (data) -- (env);
\draw[arrow] (env) -- (run);
\draw[arrow] (run) -- (out);
\draw[arrow] (out) -- (check);
\draw[arrow] (check) -- (deliver);
\end{tikzpicture}

  \caption{Il percorso di una lavorazione, dalle fotografie ai file da controllare o consegnare.}
  \label{fig:operational-workflow}
\end{figure}

La Parte I descrive l'utilizzo operativo, dalla preparazione dell'ambiente alla
consegna. La Parte II presenta gli entrypoint, il flusso applicativo e gli
oggetti Python che implementano le singole fasi.

\subsection{Terminologia}

Nel repository la parola \emph{nodo} compare con due significati. Un
\emph{nodo del dataset} è una scena completa, per esempio il nodo \code{7},
fotografata in più viste. Un \emph{nodo del grafo} è invece uno dei componenti
fisici ricostruiti all'interno di quella scena. La distinzione diventa naturale
guardando i nomi dei file:

\begin{lstlisting}[language={}]
node7_0001.jpg
node7_0002.jpg
node7_0003.jpg
\end{lstlisting}

Queste immagini formano una sola lavorazione, identificata da \code{7}. Nel
file finale possono invece comparire, per esempio, i componenti
\code{15}\#\code{1}, \code{95}\#\code{1} e
\code{116}\#\code{1}. Sono gli identificativi delle istanze
fisiche stimate dalla ReID, non nomi di immagini.

Gli ID della lavorazione possono contenere underscore. I file
\code{node15_1_0001.jpg} e \code{node15_1_0002.jpg}, per esempio, appartengono
entrambi al nodo \code{15_1}.

Nel seguito, il termine \emph{lavorazione} indica il gruppo di immagini
identificato da uno stesso ID, mentre \emph{nodo} indica normalmente un
componente del grafo.

\section{Procedura operativa}

Una normale esecuzione segue questo percorso:

\begin{enumerate}
  \item si attiva l'ambiente Conda del progetto;
  \item si indica in \filepath{.env} dove si trovano immagini, modello e
        directory di lavoro;
  \item si esegue inizialmente un solo nodo, per esempio
        \cli{python run.py 7};
  \item si apre il risultato sotto
        \filepath{<WORK_DIR>/multiview/node7/};
  \item si genera il render del grafo;
  \item si esegue la valutazione, se è disponibile il ground truth;
  \item si esporta il GeoJSON per QGIS, se richiesto.
\end{enumerate}

\subsection{Risultati attesi}

Se \envvar{WORK_DIR} conserva il valore predefinito \filepath{./output}, un run
del nodo \code{7} termina nella cartella:

\begin{center}
  \filepath{output/multiview/node7/}
\end{center}

Il risultato principale è \filepath{fused_graph.json}, che contiene componenti
e archi. \filepath{scene_nodes.json} collega le istanze fisiche alle detection
originali; \filepath{bird_view.png} consente un primo controllo della
disposizione spaziale. L'insieme degli output è descritto nel
Capitolo~\ref{chap:esecuzione}.

\manualnote[Criterio di completamento]{La sola presenza della cartella non
certifica il risultato. Per considerare conclusa una lavorazione devono
esistere \filepath{fused_graph.json} e, dopo l'esecuzione di
\cli{python visualize.py <id>}, un render leggibile. Il riepilogo finale nel
terminale deve inoltre essere coerente con il numero di viste elaborate e con
i componenti attesi o osservabili.}
